\documentclass[11pt]{article}
\usepackage{graphicx}
\usepackage{amsmath}
\usepackage{geometry}
\geometry{margin=1in}

\begin{document}
\title{OFFICIAL VERDICT: RFC \#2 --- Project Beta Full Phase Analysis}
\author{Open Physics Research Collective}
\date{\today}
\maketitle

\section{Introduction}
Project Beta investigates the transition boundary where beneficial semiconductor dopants collapse into localized carrier traps. We approach this using a unified three-phase computational framework.

\section{Phase 1: Thermal Ionization Efficiency}
We analyze the carrier contribution efficiency $\eta$ as a function of ionization energy $E_i$ at $T=300 \text{ K}$.
\begin{figure}[h]
\centering \includegraphics[width=0.55\textwidth]{research/002-doping-physics/verdicts/phase1_efficiency.png}
\caption{The rapid collapse of ionization probability beyond $0.1 \text{ eV}$.}
\end{figure}

\section{Phase 2: Wavefunction Spatial Confinement}
As the dopant potential deepens, the effective Bohr radius $a^*$ contracts. This confines the charge carrier tightly to the impurity site.
\begin{figure}[h]
\centering \includegraphics[width=0.55\textwidth]{research/002-doping-physics/verdicts/phase2_confinement.png}
\caption{Radial probability distribution comparison.}
\end{figure}

\section{Phase 3: Inter-Impurity Quantum Tunneling}
We evaluate the tunneling probability across neighboring lattice sites. For deep traps, the tunneling probability decays exponentially over typical lattice separation distances.
\begin{figure}[h]
\centering \includegraphics[width=0.55\textwidth]{research/002-doping-physics/verdicts/phase3_tunneling.png}
\caption{Tunneling suppression in deep states.}
\end{figure}

\section{Final Collective Verdict}
We formally define the boundary between dopant and trap at $E_{crit} \approx 0.1 \text{ eV}$. Beyond this boundary, spatial confinement and tunneling suppression render the impurity a recombination center rather than a carrier donor.

\end{document}
